% Options for packages loaded elsewhere
\PassOptionsToPackage{unicode}{hyperref}
\PassOptionsToPackage{hyphens}{url}
\PassOptionsToPackage{dvipsnames,svgnames,x11names}{xcolor}
\documentclass[
  10pt,
  fontset=fandol]{ctexart}
\usepackage{xcolor}
\usepackage[margin=16mm]{geometry}
\usepackage{amsmath,amssymb}
\setcounter{secnumdepth}{-\maxdimen} % remove section numbering
\usepackage{iftex}
\ifPDFTeX
  \usepackage[T1]{fontenc}
  \usepackage[utf8]{inputenc}
  \usepackage{textcomp} % provide euro and other symbols
\else % if luatex or xetex
  \usepackage{unicode-math} % this also loads fontspec
  \defaultfontfeatures{Scale=MatchLowercase}
  \defaultfontfeatures[\rmfamily]{Ligatures=TeX,Scale=1}
\fi
\usepackage{lmodern}
\ifPDFTeX\else
  % xetex/luatex font selection
\fi
% Use upquote if available, for straight quotes in verbatim environments
\IfFileExists{upquote.sty}{\usepackage{upquote}}{}
\IfFileExists{microtype.sty}{% use microtype if available
  \usepackage[]{microtype}
  \UseMicrotypeSet[protrusion]{basicmath} % disable protrusion for tt fonts
}{}
\makeatletter
\@ifundefined{KOMAClassName}{% if non-KOMA class
  \IfFileExists{parskip.sty}{%
    \usepackage{parskip}
  }{% else
    \setlength{\parindent}{0pt}
    \setlength{\parskip}{6pt plus 2pt minus 1pt}}
}{% if KOMA class
  \KOMAoptions{parskip=half}}
\makeatother
\setlength{\emergencystretch}{3em} % prevent overfull lines
\providecommand{\tightlist}{%
  \setlength{\itemsep}{0pt}\setlength{\parskip}{0pt}}
\usepackage{amsmath,amssymb,mathtools,needspace}
\xeCJKDeclareCharClass{CJK}{"2032,"0400->"04FF,"2160->"216B,"2460->"2473,"25A0,"25C7,"25CB,"25CF}
\IfFontExistsTF{Arial Unicode MS}{\xeCJKsetup{AutoFallBack=true}\setCJKfallbackfamilyfont{\CJKrmdefault}{Arial Unicode MS}}{}
\setcounter{secnumdepth}{0}
\setlength{\emergencystretch}{2em}
\allowdisplaybreaks
\usepackage{bookmark}
\IfFileExists{xurl.sty}{\usepackage{xurl}}{} % add URL line breaks if available
\urlstyle{same}
\hypersetup{
  pdftitle={多项式求值的二分法},
  colorlinks=true,
  linkcolor={Maroon},
  filecolor={Maroon},
  citecolor={Blue},
  urlcolor={Blue},
  pdfcreator={LaTeX via pandoc}}

\title{多项式求值的二分法}
\author{}
\date{}

\begin{document}
\maketitle

\subsubsection{11.2.4　多项式求值的二分法}\label{ux591aux9879ux5f0fux6c42ux503cux7684ux4e8cux5206ux6cd5}

进一步讨论多项式求值问题。仿照数列求和的做法，将所给多项式（11.2.2）的奇偶项两两合并，得

\[
P=\sum_{i=0}^{N_1-1}(a_{2i}+a_{2i+1}x)x^{2i}.
\]

若令

\[
\begin{cases}
a_i^{(1)}=a_{2i}+a_{2i+1}x,\quad i=0,1,\cdots,N_1-1,\\
x_1=x^2,
\end{cases}
\]

则有

\[
P=\sum_{i=0}^{N_1-1}a_i^{(1)}x_1^i.
\]

这样加工得出的是一个以 \(x_1\) 为变元的多项式，它与所给多项式（11.2.2）的类型相同，只是规模压缩了一半，因此上述手续是一项二分手续。

重复这种手续，二分 \(k\) 次后所给多项式被加工成

\[
P=\sum_{i=0}^{N_k-1}a_i^{(k)}x_k^i,
\]

这里

\[
\begin{cases}
a_i^{(k)}=a_{2i}^{(k-1)}+a_{2i+1}^{(k-1)}x_{k-1},\quad i=0,1,\cdots,N_k-1,\\
x_k=x_{k-1}^2.
\end{cases}
\tag{11.2.5}
\]

这样二分 \(n=\log_2 N\) 次，最终得出的系数 \(a_0^{(n)}\) 即为所求多项式的值 \(P\)。于是，多项式求值问题（11.2.2）有多项式求值的二分算法 11.3。

\textbf{算法 11.3}　对 \(k=1,2,\cdots\) 直到 \(n=\log_2 N\) 执行算式（11.2.5），结果有

\href{../../source-images/pdf-283.jpeg}{原页图像}

\[
P=a_0^{(n)}.
\]

上述算法同样可以向量化，事实上，算式（11.2.5）可表为向量化形式：

\[
\begin{pmatrix}
a_0^{(k)}\\
a_1^{(k)}\\
\vdots\\
a_{N_k-1}^{(k)}\\
x_k
\end{pmatrix}
=
\begin{pmatrix}
a_0^{(k-1)}\\
a_2^{(k-1)}\\
\vdots\\
a_{N_{k-1}-2}^{(k-1)}\\
0
\end{pmatrix}
+
\begin{pmatrix}
a_1^{(k-1)}\\
a_3^{(k-1)}\\
\vdots\\
a_{N_{k-1}-1}^{(k-1)}\\
x_{k-1}
\end{pmatrix}\cdot x_{k-1}.
\]

\end{document}
